% ============================================================
%  Admissibility All the Way Down:
%  Functional Programming and the Geometry of Reachable Programs
%
%  An interpretation of Simon Payton Jones's conversation
%  on functional programming, type systems, and useless languages
%  through the Admissibility Program framework.
%
%  Flyxion — Working Essay — lualatex
% ============================================================

\documentclass[12pt]{article}

\usepackage{fontspec}
\setmainfont{TeX Gyre Pagella}
\usepackage{microtype}
\usepackage[margin=1.3in]{geometry}
\usepackage{setspace}
\usepackage{titlesec}
\usepackage{hyperref}
\usepackage{xcolor}
\usepackage{epigraph}
\usepackage{booktabs}
\usepackage{array}
\usepackage{enumitem}

\definecolor{rsvpblue}{HTML}{1B4F72}
\definecolor{cliogreen}{HTML}{1D6A39}
\definecolor{spherered}{HTML}{7B241C}
\definecolor{repairteal}{HTML}{117A65}
\definecolor{formalismgray}{HTML}{4A4A4A}

\titleformat{\section}{\large\bfseries}{\thesection.}{0.75em}{}
\titleformat{\subsection}{\normalsize\bfseries\itshape}{\thesubsection.}{0.75em}{}

\setstretch{1.15}
\setlength{\parskip}{0.45em}
\setlength{\parindent}{1.5em}
\setlength{\epigraphwidth}{0.78\textwidth}

\hypersetup{colorlinks=true, linkcolor=rsvpblue, urlcolor=formalismgray,
  pdftitle={Admissibility All the Way Down}}

% ============================================================
\begin{document}

\begin{center}
{\LARGE\bfseries Admissibility All the Way Down}\\[0.5em]
{\large\itshape Functional Programming and the Geometry of Reachable Programs}\\[1em]
{\normalsize Flyxion \quad — \quad Working Essay}\\[0.3em]
{\small An interpretation of Simon Payton Jones, ``Functional Programming,
Thinking in Types, Useless Languages'' (2025 interview)}
\end{center}

\bigskip

\epigraph{\itshape When the limestone of imperative programming has worn away,
the granite of functional programming will be revealed underneath.}{Simon
Payton Jones}

\epigraph{\itshape Reality is Reachability. Keep the branches open.}{Flyxion,
\emph{The Admissibility Program}}

\section{Framing the Encounter}

Simon Payton Jones has spent over four decades making one bet: that programming
with values rather than mutation reveals something more fundamental about
computation than the accident of how our hardware happened to be built. In a
long 2025 interview covering functional programming, type systems, the future
of Haskell, and the implications of large language models, he offers a set of
claims that are, from the perspective of the Admissibility Program, recognisable
in unexpected depth.

This essay is not a review of the interview. It is an attempt to read Payton
Jones's account through the conceptual vocabulary of the Admissibility Program
— the research framework developed under the name Flyxion — in order to see
what each illuminates in the other. The central claim I will develop is this:
functional programming, as Payton Jones describes it, is the engineering
discipline of explicitly managing admissibility structure. Type systems are
admissibility geometries. Monads are scope-bounding mechanisms in the sense of
Spherepop. The useless-to-useful trajectory of Haskell is a worked example of
what it costs and what it reveals to take admissibility constraints seriously
over long timescales.

\section{Two Computational Primitives and Their Ontological Commitments}

Payton Jones opens with a contrast that is, at its core, a metaphysical one:
Turing machines proceed by mutation of a tape; lambda calculus proceeds by
reduction without any notion of a mutable cell. He notes — emphasising that
this is far from obvious — that both models compute exactly the same functions.
They are computationally equivalent but ontologically very different.

It is worth being precise about what kind of connection to draw between this
contrast and the Admissibility Program. The correspondence is heuristic and
structural, not formal or equivalential. The claim is not that Turing machines
\emph{are} RSVP or that lambda calculus \emph{is} Spherepop. The claim is
that the two computational models differ along the same axis that separates
RSVP from Spherepop, and that recognising this illuminates what is at stake
in both cases.

RSVP is fundamentally field-theoretic: it describes processes through
continuous variation of scalar concentration, vector flow, and entropy density
over admissibility landscapes. Its primitives are gradients, flows, and
relaxations — things that evolve continuously over a globally accessible
substrate. Spherepop is fundamentally event-theoretic: its primitives are
discrete, irreversible commitments — formation, collapse, refusal, binding,
residue. Each event eliminates possibilities and leaves a residue that
constrains what follows.

Along these dimensions, the analogy runs as follows. A Turing machine maintains
a globally accessible mutable tape; its computation is evolution over a shared
mutable substrate, tracking a continuously updated state via a program counter.
This resembles field dynamics: the tape is the substrate, the head position and
register contents are the field values, and computation is the propagation of
updates across that field. Lambda calculus, by contrast, has no global state.
Computation is constituted entirely by reduction events: each beta-reduction
irreversibly eliminates a redex, commits to a normal form, and leaves a residue
(the substituted body) that constrains subsequent reductions. There is no
persistent substrate — only an accumulating history of committed reductions.

The analogy should not be pressed to formal equivalence. A Turing machine can
be given an event-theoretic presentation; a lambda term can be given a
field-theoretic denotational semantics. What the heuristic correspondence
captures is the \emph{default orientation} of each computational primitive:
Turing's model makes global mutation its primitive operation and history a
derived record; Church's model makes irreversible commitment its primitive
operation and global state a derived convenience.

Payton Jones's lifelong wager — that programming with values reveals something
more fundamental than programming with mutation — is therefore, in Admissibility
Program terms, a claim about which orientation is primary. His hypothesis is that
the event-theoretic orientation (reduction, commitment, irreversibility) is the
bedrock, and that the mutation-based orientation (mutable state, program counter,
register modification) is a derived convenience built on top of it. Whether or
not the RSVP-Spherepop framework vindicates that hypothesis in full generality,
it provides one formal context in which the hypothesis can be stated with
precision.

\section{Side Effects as Admissibility Violations}

The core technical achievement of Haskell, as Payton Jones describes it, is to
make side effects visible in the type system. A function of type \texttt{Int ->
Int} cannot perform IO. A function of type \texttt{Int -> IO Int} can. This is
not merely a stylistic convention; it is a structural constraint on the set of
admissible programs.

In the vocabulary of the Admissibility Program, a type signature is an
admissibility specification. It defines the space of reachable futures for a
given computation: what inputs are acceptable, what outputs are possible, and —
crucially in Haskell's case — what effects are permitted along the way. A type-
incorrect program is not merely a wrong program; it is an inadmissible program,
one that has exited the set of reachable states the system is designed to
navigate.

Payton Jones identifies two consequences of this. The first is safety: 99\% of
the security exploits that arise from buffer overruns and pointer manipulation
would be eliminated by construction in a sufficiently typed language, because
those exploits are inadmissible states that a properly specified type system
simply refuses to reach. The second is maintainability: after 35 years of
development on GHC, Payton Jones still makes large-scale systematic refactorings
``fearlessly'' because the type system keeps him safe. He changes a few type
signatures and then follows a wave of type errors through the codebase,
each error pointing to a place where the system's admissibility structure has
been violated and must be restored.

This is Repair Theory in action, without the name. The type system is the repair
infrastructure of the codebase. Types define what an admissible state of the
program looks like. Type errors identify sites where admissibility has been
lost — where the program has entered an inadmissible configuration. The
programmer's task is to restore admissibility, which the compiler actively
assists by identifying every site where the repair is needed. This is
maintenance of admissible states, not construction from scratch.

Payton Jones's image of the 10-year-old codebase written by a programmer who
left long ago, with invisible coupling through shared mutable global variables,
is precisely what the Admissibility Program calls \emph{extraction}: the
projection has been mistaken for the territory, the hidden coupling has been
made invisible, and the result is that the actual admissibility constraints
of the system are no longer legible. The type system, by contrast, forces
admissibility structure into the open. The invisible coupling between functions
through shared mutable state becomes visible as a type dependency; if you change
the type, every site of the coupling lights up as an error.

\section{Scope as Geometry: Monads and Spherepop}

One of the most technically interesting moments in the interview is Payton
Jones's account of why Haskell needed monads and what they actually do. The
problem is this: functional programming eliminates side effects by default, but
real programs need IO. The question is how to admit a controlled class of
effects — a bounded region of admissible interaction with the world — while
maintaining the overall purity guarantees that make functional programs
compositionally safe.

The solution is the IO monad, which Payton Jones characterises as a type-level
marker of dirtiness. A value of type \texttt{IO Int} is a computation that,
when run, will do some IO and return an integer. The type system enforces that
these dirty computations cannot escape into pure code without explicit
acknowledgment. The do-notation sequences dirty computations in a specified
order, which is how Haskell achieves the controlled execution of effects.

This is a direct implementation of what Spherepop calls \emph{scope as
geometry}. In Spherepop, a bubble defines a topological boundary: entering the
bubble makes certain futures accessible that are not accessible from outside it,
and closing the bubble collapses those possibilities into a residue. The IO
monad creates exactly this structure at the type level. Code inside a do-block
is inside the IO bubble: it can perform effects. Code outside is outside the
bubble: it cannot. The monad's bind operation (\texttt{>>=}) is the mechanism
by which commitments made inside the bubble are threaded through subsequent
computations while keeping their effectful character visible.

The effect system that Payton Jones goes on to describe — where a type can
specify not merely \emph{whether} a computation is effectful but \emph{which
particular effects} it can perform — is a refinement of this geometry: not just
``inside or outside the IO bubble'' but a more granular topology of admissible
effectful trajectories. The library Bluefin that he mentions, and the effect
system development now happening in both Haskell and OCaml, are precisely
attempts to specify the shape of the admissibility region more precisely.

There is also a connection to what Spherepop calls \emph{refusal}. The
\texttt{unsafePerformIO} escape hatch — the function that lets you perform IO
while telling the type system to pretend you haven't — is named with explicit
acknowledgment of its danger. The name is itself a documentation of a boundary-
crossing: the programmer is asserting that they are taking on responsibility for
an action that the formal admissibility structure of the type system would refuse
to sanction. This is not quite a Spherepop refusal event in the technical sense,
but it shares the key characteristic: it is an autonomous act of stepping outside
the normal scope constraints while explicitly accepting the consequences, rather
than a silent violation of them.

A further connection worth drawing, which the interview itself only approaches
obliquely, is to the Admissibility Program's concept of Markov boundaries.
Pure code and effectful code in Haskell are separated by an informational
boundary: information can pass from pure code into effectful code freely (you
can use a pure value anywhere inside a do-block), but effectful computation
cannot flow back into pure code without explicit sanctioning through the type
system. The IO monad is, in this sense, a computational Markov boundary: it
separates two causal domains (pure and effectful) and specifies the sufficient
statistics — the type signature — that are needed to communicate across the
boundary. The monad is not merely a container for dirty computations; it is a
boundary that makes the causal structure of the program legible at the type
level. The effect systems now being developed in both Haskell and OCaml are
refinements of this boundary: instead of a binary pure/impure distinction, they
describe a more granular topology of causal domains and the specific interfaces
through which information can cross between them.

\section{Laziness, Admissibility, and Reachable Futures}

Payton Jones's account of lazy evaluation offers another point of contact. In a
lazy language, expressions are not evaluated until their values are actually
needed. The standard example he gives is a chess program that generates an
infinite tree of possible moves: in a strict language, the tree must be
generated before it can be explored, which is impossible; in a lazy language,
the generator and the explorer can be modularly separated, with the generator
producing only as much tree as the explorer actually consumes.

This is a computational implementation of the Admissibility Program's core
ontological claim: reachability is more fundamental than location. In a strict
language, all of the tree is brought into existence (evaluated, located) before
exploration begins. In a lazy language, only the reachable parts of the tree —
the parts that the explorer's trajectory actually visits — are ever instantiated.
The lazy language implements a computation in which admissibility structure
determines what exists, rather than the converse.

Payton Jones's phrase ``lazy evaluation is very powerful glue that lets you glue
together two programs that you'd like to be distinct'' is, in Admissibility
Program terms, a claim about how compositional modular design preserves the
admissibility structure of each component. Strict evaluation ``forces you to
merge them together'' — it collapses the distinction between the generator and
the explorer, destroying the modularity. This is extraction in CLIO's sense:
the projection of the combined system discards the structural independence of
its components.

\section{Type Systems as Admissibility Geometries}

Payton Jones's technical account of type systems contains, in compressed form,
a theory of what it means to specify an admissibility geometry. Let me draw out
the key structural claims.

A type system is a mechanism for specifying, at compile time, which programs are
admissible. A dynamically typed language accepts all syntactically well-formed
programs as admissible; whether a program actually terminates or produces a
sensible value is discovered only at runtime. A statically typed language
restricts the admissible programs to those that satisfy the type constraints,
eliminating at compile time a large class of programs that would fail at runtime.

The tension Payton Jones identifies — between type systems that are too weak
(rejecting programs they should admit) and type systems that are too strong
(hard to use) — is precisely the tension between an admissibility region that
is too narrow and one that is too broad. Too narrow: you exclude well-behaved
programs like ``apply list-reverse to a list of characters'' because your type
language is not expressive enough to specify that the operation is admissible
regardless of element type. Too broad: you admit programs that fail at runtime
because your admissibility specifications are insufficiently discriminating.

The goal is to find an admissibility geometry that is both maximally inclusive
of programs you want to run and maximally exclusive of programs you do not.
Parametric polymorphism (for all $\alpha$, list of $\alpha$ to list of $\alpha$)
is a move that expands the admissibility region in exactly the right direction:
it admits more programs (reverse works on lists of any type) without admitting
any bad ones (you still cannot confuse integers with floating point numbers).
Type classes and effect systems are further moves that refine the geometry.

This maps directly onto what the Admissibility Program calls the
\emph{projection-invariance problem} in CLIO: what features of an admissibility
landscape survive projection into a lower-dimensional representation, and how
can they be distinguished from artefacts of the projection operator? A type
system is precisely a projection of the full semantic landscape of a program
into a more tractable representation. The question of which type system to use
is the question of which projection preserves the most relevant admissibility
structure without introducing artefacts. Parametric polymorphism is an
admissibility-preserving projection that avoids the type-error of Pascal's
monomorphic reverse function.

\subsection{GHC Core: Admissibility Maintained Through Every Transformation}

The most striking formulation in the interview is Payton Jones's account of
GHC's internal language Core, which is a statically typed intermediate
representation of Haskell programs. The architecture is worth examining in full
because it is, without being described in these terms, a nearly complete
implementation of the Admissibility Program's concept of repair infrastructure
in a compiler.

Core is based on System F, the statically typed lambda calculus defined by
Girard and Reynolds. Every Haskell program, after parsing and type-checking,
is translated into Core. The key property is this: every optimisation pass in
GHC must transform a type-correct Core program into another type-correct Core
program. After each pass, GHC runs a type-checker on Core to verify that the
optimisation has preserved type correctness. If it has not — if the pass has
introduced a type error — that is a bug in the optimiser, caught immediately.

The consequences of this architecture are significant. An optimisation pass that
introduces a type error in Core has exited the admissible region of compiler
transformations. The type-checker for Core does not verify correctness of the
output — it cannot, in general, prove that the optimised program behaves
identically to the input. What it does is verify that the output remains
\emph{within the admissible region}: that it is still a well-typed program in
System F. This is exactly the distinction between verification (proving
correctness) and admissibility maintenance (ensuring the system has not left the
permitted region of states).

The alternative, which Payton Jones notes is universal among other production
compilers, is to have no type checker for the intermediate language at all.
In that architecture, a bug in an optimisation pass produces machine code that
may crash at runtime — often at a point far removed from the site of the
original error. The debugging process requires backtracking from the runtime
failure, through the generated machine code, through the intermediate
representations, to the specific optimisation pass that introduced the error.
This is the structure of what the Admissibility Program calls \emph{extraction}:
the admissibility constraints of the transformation have been made invisible, and
when they are violated the error propagates irreversibly into the output before
it can be detected.

With Core's type-checker, violations are detected locally and immediately. Each
optimisation pass is bracketed by an admissibility check. The error is caught at
the boundary of the pass that introduced it, not at the downstream consequence.
This is repair infrastructure in the precise sense: the type-checker is a
mechanism that continuously monitors whether the system has remained in the
admissible region, and flags violations at the earliest possible point so they
can be addressed before they propagate.

Payton Jones notes that no other production compiler has this property and
expresses genuine pride in it. From the Admissibility Program's perspective, the
rarity of this approach is itself diagnostic. We have built an enormous amount
of software infrastructure without the repair infrastructure that would allow us
to detect admissibility violations before they propagate to irreversible damage.
The internet's insecurity — which Payton Jones separately attributes to the
absence of type safety in infrastructure languages — and the absence of
type-checked intermediate representations in compilers are two faces of the same
civilisational failure: systems built without continuous admissibility monitoring,
maintained at enormous cost because violations are detected only after they have
become irreversible.

\section{The Useless-to-Useful Trajectory as Stratum Transition}

Payton Jones presents a two-dimensional design space with safety on one axis
and usefulness on the other. The first version of Haskell occupied a position
of high safety and low usefulness: a Haskell program was simply a function from
string to string, with no IO at all. C occupies a position of high usefulness
and low safety. The trajectory of Haskell's development has been to move
vertically — gaining usefulness without sacrificing safety — while Rust's
trajectory has been to move horizontally from C's position, gaining safety while
retaining usefulness.

This is a map of an admissibility landscape, though Payton Jones does not use
that language. Safety corresponds to the narrowness of the admissibility region
— to how precisely the type system specifies which programs can be written.
Usefulness corresponds to whether the admissible programs include enough of the
practically useful ones. The nirvana quadrant — safe and useful — is the region
of the admissibility landscape where the specification is simultaneously tight
enough to exclude the bad programs and generous enough to include the good ones.

The history of Haskell is the history of expanding the admissibility region in
the direction of usefulness without relaxing it in the direction of safety. Each
extension — monads for IO, type classes for polymorphism, GADTs for dependent-
like typing, the effect system now being developed — is a move that admits more
programs as safe without admitting more unsafe ones. The constraint is always
that the move must not exit the region of formal admissibility guarantees.

In Admissibility Program terms, this is a sequence of \emph{stratum
transitions}: each extension corresponds to recognising that the current
admissibility geometry was too narrow (it excluded programs that are in fact
safe) and expanding it in a way that preserves the constraints that generate the
safety guarantees. The development of the IO monad was a stratum transition: it
expanded the class of admissible Haskell programs from functions over strings to
programs with controlled effects, without relaxing the guarantee that effects
outside the IO type cannot contaminate pure code.

\section{LLMs, Admissibility Filtering, and Active Geodesic Inference}

One of the most interesting claims in the interview is Payton Jones's assertion
that static type systems are ``a huge boon for LLMs.'' His reasoning is
precise: an LLM generating code in a statically typed language can use the type
checker as a filter. It generates a candidate program, runs the compiler, and if
the program is type-incorrect it tries again. The type system dramatically
tightens the feedback cycle by rejecting inadmissible programs immediately,
before they reach any test suite. In an untyped language, the only feedback is
from running the program, which may require an elaborate test environment to
detect failures.

This is the Admissibility Program's framework of \emph{semantic navigation}
applied to code generation. The LLM is navigating a semantic manifold — the
space of possible programs — under admissibility constraints imposed by the type
system. The type checker is an admissibility filter: it immediately identifies
when a generated program has exited the admissible region and provides feedback
about where the boundary was crossed (a type error with a location). The type
errors are not obstacles to be circumvented; they are signals that guide the
navigation back into the admissible region.

This also connects directly to the framework of Active Geodesic Inference.
Intelligence is defined in that framework as the capacity to remain within a
family of dynamically admissible, low-action histories by actively reshaping the
geometry of configuration space. An LLM navigating the space of Haskell programs
with a type checker as an admissibility oracle is implementing something close to
this: at each step it generates a candidate trajectory (a candidate program) and
receives a signal about whether that trajectory is admissible. The type system
provides the geodesic structure — the metric on the program space that determines
what counts as a nearby admissible program.

Payton Jones's insight is that this loop closes much faster with a static type
system than with runtime testing. The admissibility filter is applied earlier in
the pipeline, eliminating large classes of inadmissible trajectories before they
are even committed to. This is why semantic isomers — the AGI framework's
concept of distinct internal reasoning trajectories that project to the same
external observable — matter here: two programs that produce the same output on
the test suite may have very different type structures, and the type-correct one
is more likely to remain admissible under subsequent modifications.

\section{Avoid Success at All Costs: Admissibility Under Historical Pressure}

The phrase ``avoid success at all costs'' that Payton Jones discusses is not
merely a witty aphorism. It describes a genuine tension in the management of an
admissibility region over time.

Success — in the sense of widespread adoption — creates historical pressure to
relax admissibility constraints. Users want to write programs that the current
type system rejects; the path of least resistance is to expand the admissible
region to include those programs. But if the expansion is not carefully
controlled, it admits programs that are unsafe, and the safety guarantees that
made the language valuable in the first place are eroded.

The Haskell community's commitment to principled constraint maintenance — its
refusal to simply add features because users want them, its insistence on finding
the right abstraction rather than the convenient hack — is an instance of what
the Admissibility Program calls \emph{repair theory}: the active maintenance of
an admissible state in the face of constant pressure toward inadmissibility. The
type system is the thing being repaired. The users who want unrestricted side
effects are the source of constraint erosion. The research program of finding
new type-theoretic abstractions (monads, effect systems, type classes) is the
repair mechanism: it finds ways to extend the admissible region without relaxing
the constraints that generate the safety guarantees.

Payton Jones's backward compatibility project — ensuring that a program that
compiles with GHC 10.0 also compiles with GHC 10.2 — is a different kind of
repair: preserving the admissibility of existing programs as the language
evolves. Every language extension is a potential admissibility violation for
programs that depend on the old behaviour. Maintaining backward compatibility is
the work of ensuring that the admissibility region grows only upward — admitting
new programs — without contracting downward, inadmissible programs that used to
be admissible.

This is \emph{civilisational repair} applied to programming language design.
The codebase of the world's software is the institution; the type system is the
governance mechanism; backward compatibility is the undo stack. The decision to
invest serious effort in backward compatibility is a decision to treat the
existing codebase as infrastructure worth preserving — to apply Repair Theory at
the level of an entire programming ecosystem rather than a single codebase.

\section{Types as Coordination Geometry}

One of the deepest themes in the interview is not correctness but
maintainability, and what Payton Jones says about maintainability points toward
a dimension of type systems that their usual framing as \emph{correctness tools}
does not capture. When he describes doing large-scale refactorings of GHC
``fearlessly'' because the type system keeps him safe, or when he contrasts the
difficulty of maintaining 15-year-old dynamically typed code with the relative
ease of maintaining 35-year-old statically typed code, he is not primarily
describing a correctness achievement. He is describing a coordination achievement.

A type signature is a contract. When Payton Jones writes a function with type
signature \texttt{Map k v -> k -> Maybe v}, he is not merely describing the
function's behaviour to the compiler. He is making a commitment to every future
programmer who will read, use, or modify that code — whether that programmer is
a colleague working simultaneously in a different office, a junior developer who
will join the project in five years, or Payton Jones himself in 2040. The type
signature communicates across time and organisational distance: it specifies
exactly what the function accepts, what it returns, and — through the absence of
any IO marker — that it has no side effects on any shared state.

This is what the Admissibility Program calls \emph{coordination geometry}: the
structure by which distributed agents can synchronise their local constraints
without direct communication. A type system is a propagating constraint field.
When a type changes, the type-checker identifies every site in the codebase
where that change requires a corresponding adjustment — not through runtime
discovery or test failure, but through immediate structural notification. The
type error is a coordination signal: it tells every contributor, at every level
of the stack, exactly where their local assumptions depend on the changed
constraint and exactly where repair is required.

Payton Jones's image of ``a wave of changes propagating through'' the codebase
after a type change is extraordinarily close to how Coordination Geometry
describes constraint propagation across a network of coupled agents. The type
system is the medium through which this propagation occurs. It does not merely
prevent bugs; it enables a form of coordination across time that would otherwise
require explicit communication between every pair of programmers who might ever
touch related code. In a dynamically typed codebase, that coordination is
achieved instead through documentation, convention, testing, and institutional
memory — all of which are fragile, invisible, and lossy. The type system makes
the constraint structure of the coordination explicit and formal.

This interpretation suggests a reframing of Payton Jones's most provocative
claim — that static type systems are ``a huge boon for LLMs.'' The boon is not
merely that LLMs receive faster feedback on type errors. It is that type systems
provide a formal coordination interface between human programmers and LLM-
generated code that does not exist in dynamically typed languages. An LLM
generating code with a type signature is making a legible commitment: anyone
who reads the generated code can verify, without running it, whether it connects
correctly to the surrounding code's constraint structure. The type system is the
medium through which human programmers can coordinate with an LLM tool across
the boundary of their respective competencies.

Viewed at civilisational scale: the reason functional programming ideas —
lambdas, polymorphism, monads, garbage collection — have steadily infected
mainstream imperative languages is not that they are theoretically elegant. It
is that they reduce coordination costs. Each of those ideas makes some aspect
of a program's constraint structure more visible and more formal, allowing
programmers to coordinate more effectively across time, distance, and
organisational boundaries. The ``limestone'' of imperative programming wears away
not because functional programming is aesthetically superior, but because the
constraint-visibility offered by functional ideas reduces the friction of
civilisational-scale software maintenance.

\section{When the Formalism Becomes the Interface}

A recurring theme in the interview, stated most clearly in Payton Jones's
account of how he designs programs, is that experienced Haskell programmers
stop using types as annotations and start thinking through types. He writes the
types first. He changes a type and follows the wave of errors. He describes the
type system as his ``design language.'' This is not a description of a tool for
checking designs already conceived; it is a description of a medium in which
design itself happens.

This is the thesis of the Admissibility Program's essay \emph{When the Formalism
Becomes the Interface}: there is a phase transition in the relationship between
a formal system and its user, after which the formalism ceases to be a
representation of thought and becomes the medium through which thought occurs.
The formalism does not describe the interface; it \emph{is} the interface.

For Haskell programmers at that phase transition, the admissibility geometry of
the type system is not an external constraint imposed on their thinking. It is
the internal structure through which their thinking is organised. A type error is
not ``my program failed a check''; it is ``I have not yet found the right
decomposition of this problem into admissible pieces.'' The type system becomes
an active collaborator in the process of design — something that Payton Jones
gestures at when he says he finds that ``changing the representation of a data
structure'' and then compiling tells him exactly where in the codebase that
representation is used, written, and freshly allocated.

Null Convention Logic, which Payton Jones cites in the context of hardware for
functional programming, is another instance of this phenomenon. In NCL, a
circuit is complete — its computation is abstracted — at the moment of signal
stabilisation: the formalism of completion-governed logic becomes the medium
through which the circuit expresses its computational boundary. The Admissibility
Program's \emph{Abstraction as Reduction} monograph identifies NCL stabilisation
as one of five formally distinct instances of the same invariant: inner
complexity is resolved to allow outer complexity to grow without conflict. When
Haskell's type system becomes the programmer's design language, that same
invariant is at work: the inner complexity of type resolution is completed by
the compiler so that the outer complexity of system design can proceed without
interference.

The phase transition that makes a formalism into an interface is not available to
all formalisms. It requires that the formalism be expressive enough to capture
the relevant admissibility structure, and complete enough that working within it
feels like thinking rather than translating. Payton Jones's assessment that
Haskell embodies the single core principle of ``no unrestricted side effects'' and
declines to compromise it is, from this perspective, a prerequisite for the phase
transition to occur: a formalism that makes exceptions for convenience cannot
become the medium of thought, because thinking through it would require constantly
tracking which parts of the formalism apply and which have been suspended.

\section{What Remains Unreachable: The Boundaries of the Admissibility Program's Reach}

Payton Jones's account also contains points of genuine tension with the
Admissibility Program's framework, worth naming explicitly.

The first is performance. He acknowledges that lazy evaluation, despite its
theoretical elegance, creates practical problems for performance prediction and
memory management. The admissibility constraints that make Haskell's programs
compositionally safe also make their performance behaviour less predictable,
because the order of evaluation is determined not by the program's structure but
by what the consumer actually demands. This is a case where the admissibility
geometry of the type system is genuinely constraining in a way that is not
straightforwardly captured by the Admissibility Program's framework, which
tends to treat constraints as revealing structure rather than hiding it.

The second is the question of what the type system cannot specify. Payton Jones
acknowledges that no type system can prevent a program that deliberately
exfiltrates its entire database in response to a request. Deadlock, high-level
logic errors, and semantic bugs are outside the reach of any type-based
admissibility specification. The Admissibility Program's notion of admissibility
is also formal — it specifies what transitions are allowed from a given state —
but the relationship between formal admissibility and actual correctness or
safety is not straightforwardly compositional. You can have a formally admissible
trajectory through a semantic space that leads somewhere genuinely terrible.

The third is the gap between the theory and the ecosystem. Payton Jones is frank
that Haskell is talked about far more than it is used. The principled constraint
maintenance that makes it theoretically attractive makes it practically
difficult, and the community that finds it rewarding is self-selected for
people who want to think in a different way. The Admissibility Program is
similarly a framework that requires a substantial investment in new ontological
vocabulary before it becomes productive. Neither has solved the problem of making
the insight accessible without diluting it.

\section{The Limestone and the Granite}

Payton Jones's most evocative image — ``when the limestone of imperative
programming has worn away, the granite of functional programming will be
revealed underneath'' — is a geological metaphor that the Admissibility Program
has independently arrived at. The five-layer stratigraphy of the Admissibility
Program places practice and embodied knowledge at the bedrock, core intuitions
as subsoil, primary theories as mantle, formalization programs as crust, and
domain applications as surface. The claim is that the deeper layers are more
durable, more structurally fundamental, and more resistant to erosion than the
surface features.

Payton Jones's geological intuition is the same: imperative programming is
limestone — useful, widely distributed, built from the accumulated sediment of
decades of engineering practice, but ultimately soluble. Lambda calculus, value-
based computation, and the mathematical structures of functional programming are
granite: older, denser, formed under conditions of higher pressure, and
fundamentally more resistant to the acids of time and use.

The Admissibility Program's claim is that admissibility structure is the
deepest geological layer of all. Not lambda calculus specifically, not
functional programming as a programming paradigm, but the geometry of what
can be reached from where — the structure of admissible transformations across
physics, cognition, memory, computation, and civilisation. Lambda calculus is
one particularly clean way of making that structure visible in the domain of
computation. Payton Jones has spent his career ensuring that the formal
admissibility structure of programs is not hidden in a mutable global state
somewhere but made visible, legible, and repairable through the type system.

That is not a narrow contribution to programming language design. It is a
contribution to the broader project of making admissibility structure legible
wherever it lives.

\bigskip
\noindent\textcolor{formalismgray}{\rule{\linewidth}{0.4pt}}
\smallskip
\noindent{\small\textit{This essay interprets Simon Payton Jones, ``Co-Creator
of Haskell: Functional Programming, Thinking in Types, Useless Languages,''
2025 interview (transcript). All quotations are from that transcript. The
Admissibility Program frameworks referenced are documented in \emph{Project Map
and Intellectual Architecture of the Flyxion Research Program} (Flyxion, 2026).}}

\end{document}

